\section{El producto de números enteros}

La interpretación del producto como suma repetida nos permitió construir el producto de números naturales. Sin embargo, ya no basta para explicar expresiones como
\[
(-3)\cdot4,
\qquad
3\cdot(-4)
\qquad\text{o}\qquad
(-3)\cdot(-4).
\]
No tiene sentido interpretar literalmente la primera de ellas como ``sumar \(4\) menos tres veces''. Necesitamos extender el producto a los números enteros.

Al hacerlo, no abandonaremos las propiedades que ya conocemos: el producto seguirá siendo asociativo y conmutativo, tendrá al \(1\) como elemento neutro, se distribuirá respecto de la suma y continuará cumpliendo que el producto por cero es cero. Estas propiedades permitirán determinar qué debe significar multiplicar por un número negativo.

Así, el producto puede considerarse ahora como una operación
\[
\cdot\colon\mathbb{Z}\times\mathbb{Z}\to\mathbb{Z}.
\]
En particular, el producto de dos enteros vuelve a ser un entero.

\subsection{El opuesto como producto por \texorpdfstring{\(-1\)}{-1}}

Recordemos que \(-1\) es el opuesto de \(1\). Por lo tanto,
\[
(-1)+1=0
\]
Multipliquemos ambos sumandos por un entero \(a\). Por la propiedad distributiva,
\begin{gather*}
a\cdot\bigl((-1)+1\bigr)
=a\cdot(-1)+a\cdot1
\end{gather*}
Operando usando las propiedades, obtenemos
\begin{align*}
  a\cdot\bigl((-1)+1\bigr) &= a\cdot(-1)+a\cdot 1 \\ 
  a \cdot (0) &= a\cdot(-1) + a\cdot1 \\ 
  0 &= a\cdot(-1) + a \\ 
  -a + 0 &= a\cdot(-1) +a - a \\ 
  -a &= a\cdot(-1)
\end{align*}
Es decir, \(a\cdot(-1)\) es el opuesto de \(a\). Por la definición de opuesto,
\[
\boxed{a\cdot(-1)=-a}
\]
Esta igualdad conecta el producto con la idea estudiada en el capítulo anterior: un entero negativo puede escribirse como el producto de su opuesto positivo por \(-1\).

\subsection{Producto por un opuesto}

Si \(b\) es un entero, acabamos de ver que
\[
-b=b\cdot(-1)
\]
Entonces, usando la propiedad asociativa,
\begin{align*}
a\cdot(-b)
&=a\cdot\bigl(b\cdot(-1)\bigr)\\
&=(a\cdot b)\cdot(-1)\\
&=-(a\cdot b)
\end{align*}
Por la propiedad conmutativa se obtiene también
\[
(-a)\cdot b=-(a\cdot b)
\]
Por lo tanto, multiplicar exactamente uno de los factores por su opuesto produce el opuesto del producto.

\subsection{Producto de dos opuestos}

Ahora podemos determinar el producto de dos números negativos sin memorizar una regla nueva. Como \(-a=a\cdot(-1)\), tenemos
\begin{align*}
(-a)\cdot(-b)
&=\bigl(a\cdot(-1)\bigr)\cdot(-b)\\
&=(-1)\cdot\bigl(a\cdot(-b)\bigr)
\end{align*}
En la segunda igualdad usamos las propiedades asociativa y conmutativa. Por la sección anterior,
\[
a\cdot(-b)=-(a\cdot b)
\]
Por consiguiente,
\begin{align*}
(-a)\cdot(-b)
&=(-1)\cdot\bigl(-(a\cdot b)\bigr)\\
&=-\bigl(-(a\cdot b)\bigr)\\
&=a\cdot b
\end{align*}
Así, el producto de dos opuestos es el producto de los números originales:
\[
\boxed{(-a)\cdot(-b)=a\cdot b}
\]

Las cuatro posibilidades para los signos se resumen ahora en consecuencias de los opuestos:
\begin{gather*}
(+a)\cdot(+b)=a\cdot b\\
(-a)\cdot(+b)=-(a\cdot b)\\
(+a)\cdot(-b)=-(a\cdot b)\\
(-a)\cdot(-b)=a\cdot b
\end{gather*}
Por ejemplo,
\[
(-3)\cdot(-4)=3\cdot4=12
\]
